\documentclass[11pt]{article}

\usepackage[a4paper,margin=1in]{geometry}
\usepackage{amsmath,amssymb}
\usepackage{enumitem}
\usepackage{fancyhdr}
\usepackage{xcolor}
\usepackage{listings}
\lstset{
  basicstyle=\ttfamily,
  columns=fullflexible,
  keepspaces=true,
  mathescape=true,
  showstringspaces=false,
  frame=single
}

\pagestyle{fancy}
\fancyhf{}
\lhead{CSE 400: Fundamentals of Probability in Computing}
\rhead{Lecture Scribe - 10}
\cfoot{\thepage}

% ===================== CUSTOM COMMANDS =====================
\newcommand{\solution}{
    \vspace{0.5cm}
    \noindent\textbf{\textcolor{blue}{Solution:}}
    \vspace{0.2cm}
    
}

% ===================== TITLE =====================
\title{
    \normalsize School of Engineering and Applied Science (SEAS), Ahmedabad University \\
    \vspace{0.2cm}
    \textbf{CSE 400: Fundamentals of Probability in Computing}\\
    \Large Homework Assignment-1 Submission
}
\author{} 
\date{}

\begin{document}
\maketitle

\vspace{-2cm}
\begin{center}
    \begin{tabular}{ll}
        \textbf{Name:} & {Krishna Patel \hspace{2.5in}} \\ [1.5ex]
        \textbf{Enrollment No:} & {AU2440032 \hspace{2.5in}} \\ [1.5ex]
        \textbf{Email:} & {krishna.p@ahduni.edu.in \hspace{2.1in}} \\ [1.5ex]
        \textbf{Date of Submission:} & {Feburary 12th\textsuperscript{th}, 2026 (11:59 PM) \hspace{1.2in}}
    \end{tabular}
\end{center}

\hrule
\vspace{0.5cm}
\begin{center}
{\LARGE\bfseries  Randomized Min-Cut Algorithm}
\end{center}

\bigskip
\section{Introduction}

This lecture covers the \textbf{min-cut problem} and algorithms used to solve it.
The lecture includes:

\begin{itemize}
    \item Definition and applications of the min-cut problem
    \item Successful and unsuccessful min-cut runs
    \item The \textbf{Max-Flow Min-Cut Theorem}
    \item Deterministic min-cut algorithm (Stoer–Wagner)
    \item Randomized min-cut algorithm (Karger’s algorithm)
    \item Pseudocode for deterministic and randomized approaches
    \item Comparison between deterministic and randomized min-cut algorithms
    \item Theorem related to probability of obtaining a min-cut
    \item Reference to Python simulation activity
\end{itemize}

\section{Core Definitions and Concepts}

\subsection{Min-Cut Problem: Why use min-cut?}

Min-cut algorithms are used in various applications related to:

\begin{itemize}
    \item \textbf{Network connectivity}
    \item \textbf{Reliability}
    \item \textbf{Optimization}
\end{itemize}

Applications explicitly mentioned:

\begin{itemize}
    \item \textbf{Network Design:}  
    Min-cut helps in improving the efficiency of communication and optimizing network flow. The algorithm is used in network design to find the minimum capacity cut.
    \item \textbf{Communication Networks:}  
    For understanding the vulnerability of networks to failures, minimum cut can be useful. It helps in building robust and fault-tolerant communication networks.
    \item \textbf{VLSI Design:}  
    In Very Large Scale Integration (VLSI) design, the algorithm is useful for partitioning circuits into smaller components, leading to reduced interconnectivity complexity.
\end{itemize}

(Reference mentioned in lecture: Section 1.5, \textit{Application: A Randomized Min-Cut Algorithm}, \textit{Probability and Computing}, 2nd Edition.)

\subsection{What is Min-Cut?}

\begin{itemize}
    \item A \textbf{cut-set} in a graph is a set of edges whose removal breaks the graph into two or more connected components.
    \item Given a graph \( G = (V, E) \) with \( n \) vertices, the \textbf{minimum cut (min-cut) problem} is to find a minimum cardinality cut-set in \( G \).
\end{itemize}

Additional lecture statement:

\begin{itemize}
    \item Min-cut algorithms, such as \textbf{Karger’s algorithm}, are random and can be sensitive to the initial choice of edges. If the algorithm contracts critical edges early in the process, it might find a smaller cut.
\end{itemize}

\subsection{Edge Contraction}

\begin{itemize}
    \item The main operation in the algorithm is \textbf{edge contraction}.
    \item Edge contraction is an operation that removes an edge from a graph while simultaneously merging the two vertices.
\end{itemize}

In contracting an edge \( (u, v) \):

\begin{itemize}
    \item Merge vertices \( u \) and \( v \) into one vertex
    \item Eliminate all edges connecting \( u \) and \( v \)
    \item Retain all other edges in the graph
\end{itemize}

Resulting graph properties:

\begin{itemize}
    \item The new graph \textbf{may have parallel edges}
    \item The new graph has \textbf{no self-loops}
\end{itemize}

\section{Main Results / Theorems}

\subsection{Successful Min-Cut Run}

\begin{itemize}
    \item A \textbf{successful min-cut run} refers to the success in the outcome of an algorithm designed to find the minimum cut in a graph.
\end{itemize}

(Figure labeled: \textit{Successful Run} shown in lecture.)

\subsection{Unsuccessful Min-Cut Run}

\begin{itemize}
    \item An \textbf{unsuccessful min-cut run} refers to an iteration of a min-cut algorithm where the algorithm \textbf{fails to correctly identify the minimum cut} of a given graph.
\end{itemize}

(Figure labeled: \textit{Unsuccessful Run} shown in lecture.)

\subsection{Max-Flow Min-Cut Theorem}

\textbf{Statement:}

\begin{quote}
``In a flow network, the maximum amount of flow passing from the source to the sink is equal to the total weight of the edges in a minimum cut.''
\end{quote}

Definitions provided:

\begin{itemize}
    \item \textbf{Capacity of a cut:}  
    The sum of the capacity of the edges in the cut that are oriented from a vertex \( \in X \) to a vertex \( \in Y \)
    \item \textbf{Minimum cut:}  
    The cut in the network that has the smallest possible capacity
    \item \textbf{Minimum cut capacity:}  
    The capacity of the minimum cut
    \item \textbf{Maximum flow:}  
    The largest possible flow from source \( S \) to sink \( T \)
\end{itemize}

\section{Proofs / Derivations}

\subsection{Stoer--Wagner Min-Cut Algorithm (Theorem Statement)}

Let \( s \) and \( t \) be two vertices of a graph \( G \).  
Let \( G / \{s, t\} \) be the graph obtained by merging \( s \) and \( t \).

\textbf{Result:}

\begin{itemize}
    \item A minimum cut of \( G \) can be obtained by taking the smaller of:
    \begin{itemize}
        \item A minimum \( s \)--\( t \) cut of \( G \), and
        \item A minimum cut of \( G / \{s, t\} \)
    \end{itemize}
\end{itemize}

\textbf{Justification (as stated):}

\begin{itemize}
    \item Either there is a minimum cut of \( G \) that separates \( s \) and \( t \), then a minimum \( s \)--\( t \) cut of \( G \) is a minimum cut of \( G \);
    \item Or there is none, then a minimum cut of \( G / \{s, t\} \) does the job.
\end{itemize}

\section{Algorithms / Pseudocode}

\subsection{Deterministic Min-Cut Algorithm}

\textbf{Algorithm 1: MinimumCutPhase(\( G, a \))}

\begin{lstlisting}
    A $\leftarrow$ {a};
while A $\neq$ V do
    add to A the most tightly connected vertex;
return the cut weight as the "cut of the phase";
\end{lstlisting}



\textbf{Algorithm 2: MinimumCut(\( G \))}

\begin{lstlisting}
    while $|V| \geq 1$ do
    choose any a from V;
    MinimumCutPhase(G, a);
    if the cut-of-the-phase is lighter than the current minimum cut then
        store the cut-of-the-phase as the current minimum cut;
    shrink G by merging the two vertices added last;
return the minimum cut;
\end{lstlisting}


\subsection{Randomized Min-Cut Algorithm}

\textbf{Why Randomized Algorithm?}

Statements as given:

\begin{itemize}
    \item Randomized algorithms provide a probabilistic guarantee of success.
    \item It provides a more accurate estimate of the minimum cut with fewer iterations.
\end{itemize}

Additional points listed:

\begin{itemize}
    \item Efficiency
    \item Parallelization
    \item Approximation Guarantees
    \item Avoidance of Worst-Case Instances
    \item Heuristic Nature
    \item Robustness
\end{itemize}

\textbf{Karger’s Randomized Algorithm}

(Example run illustrated in lecture slides showing random edge selection and contraction. Caption indicates that the output cut may not be minimal.)

\textbf{Algorithm 3: Recursive-Randomized-Min-Cut(\( G, \alpha \))}

\textbf{Input:}  
An undirected multigraph \( G \) with \( n \) vertices, and an integer constant \( \alpha > 0 \)

\textbf{Output:}  
A cut \( C \) of \( G \)


\begin{lstlisting}
if n <= $\alpha^3$ then
    C <- a min-cut of G found using brute force (exhaustive) search;
else
    for i <- 1 to $\alpha$ do
        G' <- multigraph obtained by applying
        n - $\lceil n / \sqrt{\alpha} \rceil$ random contraction steps on G;
        C' <- Recursive-Randomized-Min-Cut(G', $\alpha$);
        if i = 1 or |C'| < |C| then
            C <- C';
return C;
\end{lstlisting}





\section{Worked Examples}

\begin{itemize}
    \item Figures illustrating:
    \begin{itemize}
        \item Successful min-cut run
        \item Unsuccessful min-cut run
        \item Example run of Karger’s Randomized Algorithm where the output cut is not minimal
    \end{itemize}
\end{itemize}

\section{Conclusion}

\subsection{Comparison: Deterministic vs Randomized Min-Cut}

General statement:

\begin{itemize}
    \item Any specific problem is responsible for deciding the approach.
\end{itemize}

\textbf{Deterministic Min-Cut:}

\begin{itemize}
    \item Always guarantees an exact minimum cut
    \item May have higher time complexity for large graphs
    \item Stoer--Wagner algorithm time complexity:
    \( O(V \cdot E + V^2 \log V) \)
\end{itemize}

\textbf{Randomized Min-Cut:}

\begin{itemize}
    \item Provides an approximate minimum cut with high probability
    \item Karger’s algorithm time complexity:
    \( O(V^2) \)
\end{itemize}

\subsection{Theorem for Min-Cut Set}

\begin{itemize}
    \item The algorithm outputs a min-cut set with probability at least:
\end{itemize}

\[
\frac{2}{n(n - 1)}
\]

\subsection{Python Simulation}

\begin{itemize}
    \item Students are instructed to:
    \begin{itemize}
        \item Open the Campuswire post regarding Lecture 10
        \item Download the \texttt{.ipynb} file pasted in the comments
    \end{itemize}
\end{itemize}

\bigskip
\hrule
\vspace{0.2cm}
\begin{center}
    \small \textit{End of Submission}
\end{center}

\end{document}
